\chapter{LITERATURE REVIEW AND THEORETICAL BACKGROUND}

\section{Literature Review}
\subsection{Traditional and Smart Prepaid Metering}
Historically, electrical utility grids have relied on electromechanical induction meters that require manual reading and postpaid billing cycles. This conventional approach places the financial risk entirely on the utility provider, leading to revenue collection challenges, delayed payments, and high debt-recovery costs. To mitigate these risks, the industry introduced smart prepaid metering. In these systems, consumers purchase energy credits in advance, which are digitally loaded onto the meter. While prepaid metering successfully reduces consumer defaults, the traditional digital implementation remains heavily centralized. The balance calculations and tariff deductions are processed on private, closed-source corporate servers, leaving consumers with no ability to independently verify the accuracy of their consumption deductions. Recent smart-metering implementations demonstrate the use of IoT telemetry and automated billing for this purpose 

\subsection{IoT-Based Energy Monitoring}
The integration of the Internet of Things (IoT) has driven the transition from basic smart meters to Advanced Metering Infrastructure (AMI). IoT-based energy monitoring utilizes microcontrollers and edge sensors to continuously sample analog electrical parameters and transmit them wirelessly to central hubs. Literature shows that utilizing devices like the ESP32 combined with analog-to-digital converters allows for highly accurate, low-latency tracking of live power demand \cite{zhao2020energy}. However, a significant vulnerability persists in modern AMI networks: they act merely as data pipelines. The vast amounts of sensitive consumption telemetry they generate are still funneled into centralized databases. This centralization not only creates a single point of failure susceptible to cyber-attacks but also preserves the "black-box" nature of utility billing.

\subsection{Blockchain for Decentralized Energy Systems}
To solve the trust and security deficits inherent in centralized AMI, recent research has shifted toward decentralized ledgers and Web3 architectures. Blockchain technology provides an immutable, distributed, and peer-to-peer network to ensure security and transparency among grid users. In utility applications, self-executing smart contracts replace central billing servers. These deterministic programs automatically recalculate a consumer's token balance based on incoming telemetry and predetermined tariff rules. Because the ledger is public and cryptographically secured, neither the utility provider nor the consumer can alter historical consumption data or manipulate the billing mathematics, effectively establishing a trustless utility framework \cite{mololoth2023blockchain,dong2023blockchain}.

\subsection{Research Gaps}
While IoT telemetry, blockchain billing, and machine learning anomaly detection have been extensively researched, they are predominantly studied in isolation. This has left several critical gaps in modern smart grid architecture:

\begin{itemize}
    \item The Physical Enforcement Gap: Pure Web3 and blockchain architectures handle financial state changes perfectly but lack the physical control mechanisms necessary to enforce those states. A smart contract can reduce a balance to zero, but it cannot physically cut the power without a dedicated hardware integration layer.
    \item The Localization Gap: Traditional anomaly detection systems often rely on cloud-based, centralized processing with high latency, which delays the identification of localized power theft \cite{jithish2023anomaly,husnoo2023misbehaviour}.
    \item The Trust-Automation Gap: Existing IoT meters automate the reading process but fail to automate the trust process, still relying on private databases for final tariff execution.
\end{itemize}


\section{Theoretical Background}
\subsection{Internet of Things Edge Sensing}
The Internet of Things (IoT) extends internet connectivity beyond standard devices to any range of traditionally non-internet-enabled physical devices. In smart grid applications, edge sensing involves deploying microcontrollers, such as the ESP32, directly at the consumer's load source. These edge nodes utilize Analog-to-Digital Converters (ADCs) to sample continuous analog signals—such as voltage or current variations—and digitize them into discrete time-series data. Beyond passive data collection, intelligent edge sensing incorporates actuation capabilities. By integrating electromechanical relays, the microcontroller can physically break or close an electrical circuit based on digital commands, effectively serving as the physical enforcement layer for automated grid management.

A primary advantage of edge sensing over traditional centralized telemetry is the utilization of localized data processing. Instead of transmitting heavy, raw analog arrays over a wireless network, the microcontroller performs lightweight preprocessing directly at the data source. The edge node aggregates the instantaneous readings, calculates the cumulative energy consumption ($\Delta E$), and packages these processed metrics into lightweight JSON payloads. This localized distillation minimizes bandwidth consumption and ensures low-latency transmissions to the middleware server.

\subsection{RESTful Backend Architecture}
To process and route the continuous stream of telemetry generated by the IoT edge, a robust middleware layer is required. Representational State Transfer (REST) is a software architectural style that defines a set of constraints for creating web services. A RESTful backend, implemented using frameworks like FastAPI, operates statelessly, meaning every HTTP request from the edge device contains all the information necessary to execute the request. This middleware acts as the central nervous system of the smart grid, asynchronously ingesting JSON payloads from the ESP32, passing data arrays to the machine learning engine for real-time analysis, and communicating with the blockchain network via Remote Procedure Calls (RPC).[6]

\subsection{Blockchain and Smart Contracts}
A blockchain is a decentralized, cryptographically secure distributed ledger. Rather than relying on a centralized database administrator, transactions are validated by the network consensus. A smart contract is self-executing deterministic code deployed on this ledger. In the context of a prepaid utility grid, a Solidity-based smart contract operates as an immutable state machine. It stores a user's cryptographic token balance and recalculates it dynamically based on incoming energy consumption telemetry. The automated tariff deduction logic follows the core mathematical evaluation:
\[
Bal_{new} = Bal_{current} - (\Delta E \times R_{fees})
\]
Where $\Delta E$ represents the cumulative energy consumed over a specific epoch and $R_{fees}$ represents the predefined tariff rate. Because this computation occurs on a distributed ledger, the billing process is completely transparent and mathematically guaranteed.

\subsection{Machine Learning Components}
Modern grid intelligence relies on applying mathematical models to incoming telemetry to automate decision-making.

\subsubsection{Current-Differential Anomaly Detection}
According to Kirchhoff's Current Law, in a standard closed-loop, single-phase electrical system without faults, the current flowing into the load through the active line wire must equal the current returning through the neutral wire. To detect energy theft—specifically meter bypassing or line hooking—this project employs a threshold-based current-differential anomaly detector, an approach also examined in smart-meter research \cite{ayanlade2021energytheft}. The algorithm continuously evaluates the absolute difference between the line current ($I_{line}$) and the neutral current ($I_{neutral}$). A tampering anomaly is identified if this difference exceeds a predetermined natural leakage tolerance:
\[
\left\lvert I_{line} - I_{neutral} \right\rvert > \text{threshold}
\]
This deterministic evaluation serves as the primary trigger for autonomous theft alerts on the administrative dashboard.

\subsubsection{Linear Regression for Consumption Forecasting}
Linear Regression is a fundamental statistical approach used to model the relationship between a scalar response and one or more explanatory variables. In this system, it is utilized as a supervised forecasting model. By processing historical time-series data of a consumer's power usage, the algorithm maps a continuous consumption trajectory. Research on electrical-load forecasting demonstrates the value of learning from historical consumption patterns for scalable prediction \cite{gholizadeh2022federated,zheng2024federated}. It estimates the temporal horizon when the dependent variable (the token balance) may intersect with zero. This regression analysis transforms raw telemetry into an actionable "Estimated Time Remaining" metric, improving the prepaid consumer experience.

